\clearpage
\SourceLeaf{105}%

\renewcommand{\BookRunningTitle}{Theory of Biquadratic Residues}
\renewcommand{\Article}[1]{\par\medskip\Needspace{4\baselineskip}\pdfbookmark[1]{Article #1}{paper-five-art-#1}\centerline{\textbf{#1.}}\nobreak\smallskip}
\begin{center}
\textbf{THEORY OF BIQUADRATIC RESIDUES.}\par
\bigskip SECOND PAPER.
\end{center}
\BookRule{.16\linewidth}
\Article{24}
In the first paper we completely dealt with everything required for the biquadratic classification of the number $+2$. Namely, we think of all numbers not divisible by the modulus $p$ (assumed to be a prime number of the form $4n+1$) as distributed among four collections $A,B,C,D$, according as each, raised to the power with exponent $\frac14(p-1)$, becomes congruent modulo $p$ to $+1$, $+f$, $-1$, $-f$, with $f$ denoting either root of the congruence $ff\equiv-1\ (\mathrm{mod.}\ p)$. We found that the decision as to which collection the number $+2$ belongs to depends on splitting the number $p$ into two squares: if we put $p=aa+bb$, with $aa$ denoting the odd square and $bb$ the even one, and further assume that the \emph{signs} of $a,b$ are chosen so that $a\equiv1\ (\mathrm{mod.}\ 4)$, $b\equiv af\ (\mathrm{mod.}\ p)$, then the number $+2$ must belong to the collection $A,B,C,D$ according as $\frac12b$ has the form $4n$, $4n+1$, $4n+2$, $4n+3$, respectively.

A rule for classifying the number $-2$ also follows immediately. Since $-1$ belongs to class $A$ when $\frac12b$ is even, and to class $C$ when it is odd, the theorem in Article 7 shows that the number $-2$ belongs to class $A,B,C,D$ according as $\frac12b$ has the form $4n$, $4n+3$, $4n+2$, $4n+1$, respectively.

These theorems can also be expressed as follows:

\SourceLeaf{106}%
\begin{center}
\begin{tabular}{c|c|c}
Belongs&$+2$&$-2$\\\hline
to the collection&\multicolumn{2}{c}{if $b$, modulo $8$, is congruent to}\\\hline
$A$&$0$&$0$\\
$B$&$2a$&$6a$\\
$C$&$4a$&$4a$\\
$D$&$6a$&$2a$
\end{tabular}
\end{center}

We readily see that, stated in this way, the theorems no longer depend on the condition $a\equiv1\ (\mathrm{mod.}\ 4)$, but still hold if $a\equiv3\ (\mathrm{mod.}\ 4)$, provided that the other condition, $af\equiv b\ (\mathrm{mod.}\ p)$, is maintained.

It is equally easy to see that the content of these theorems can be elegantly condensed into a single formula, namely:
\begin{quote}\itshape if $a$ and $b$ are taken positive, we always have\end{quote}
\[
b^{\frac12ab}\equiv a^{\frac12ab}\,2^{\frac14(p-1)}\ (\mathrm{mod.}\ p)
\]

\Article{25}
Let us now see how far inductive investigation points towards the classification of the number $3$. The table in Article 11, continued further (always taking the least primitive root), shows that $+3$ belongs
\begin{center}to the collection\end{center}
\[
\begin{array}{r|r|r||r|r|r||r|r|r||r|r|r}
\multicolumn{3}{c||}{A\ \text{for}}&\multicolumn{3}{c||}{B\ \text{for}}&\multicolumn{3}{c||}{C\ \text{for}}&\multicolumn{3}{c}{D\ \text{for}}\\
p&a&b&p&a&b&p&a&b&p&a&b\\
13&-3&+2&17&+1&-4&37&+1&-6&5&+1&+2\\
109&-3&+10&29&+5&+2&61&+5&-6&41&+5&-4\\
181&+9&+10&53&-7&+2&73&-3&-8&149&-7&+10\\
193&-7&-12&89&+5&-8&97&+9&+4&173&+13&+2\\
229&-15&+2&101&+1&+10&157&-11&-6&&&\\
277&+9&+14&113&-7&-8&241&-15&-4&&&\\
&&&137&-11&-4&&&&&&\\
&&&197&+1&-14&&&&&&\\
&&&233&+13&+8&&&&&&\\
&&&257&+1&-16&&&&&&\\
&&&269&+13&+10&&&&&&\\
&&&281&+5&+16&&&&&&\\
&&&293&+17&+2&&&&&&
\end{array}
\]

\SourceLeaf{107}%
At least at first sight, we see no simple connection between the values of the numbers $a,b$ to which the same collection corresponds. But, considering that a similar decision in the theory of quadratic residues is made by a simpler rule for the number $-3$ than for the number $+3$, we may hope for equally favourable success in the theory of biquadratic residues. We find that $-3$ belongs to the collection
\[
\begin{array}{r|r|r||r|r|r||r|r|r||r|r|r}
\multicolumn{3}{c||}{A\ \text{for}}&\multicolumn{3}{c||}{B\ \text{for}}&\multicolumn{3}{c||}{C\ \text{for}}&\multicolumn{3}{c}{D\ \text{for}}\\
p&a&b&p&a&b&p&a&b&p&a&b\\
37&+1&-6&5&+1&+2&13&-3&+2&29&+5&+2\\
61&+5&-6&17&+1&-4&73&-3&-8&41&+5&-4\\
157&-11&-6&89&+5&-8&97&+9&+4&53&-7&+2\\
193&-7&-12&113&-7&-8&109&-3&+10&101&+1&+10\\
&&&137&-11&-4&181&+9&+10&197&+1&-14\\
&&&149&-7&+10&229&-15&+2&269&+13&+10\\
&&&173&+13&+2&241&-15&-4&293&+17&+2\\
&&&233&+13&+8&277&+9&+14&&&\\
&&&257&+1&-16&&&&&&\\
&&&281&+5&+16&&&&&&
\end{array}
\]
where the law suggested by inductive investigation presents itself at once. Namely, $-3$ belongs to the collection
\[
\begin{array}{ll}
A,&\text{whenever }b\text{ is divisible by }3\text{, equivalently }b\equiv0\ (\mathrm{mod.}\ 3)\\
B,&\text{whenever }a+b\text{ is divisible by }3\text{, equivalently }b\equiv2a\ (\mathrm{mod.}\ 3)\\
C,&\text{whenever }a\text{ is divisible by }3\text{, equivalently }a\equiv0\ (\mathrm{mod.}\ 3)\\
D,&\text{whenever }a-b\text{ is divisible by }3\text{, equivalently }b\equiv a\ (\mathrm{mod.}\ 3)
\end{array}
\]

\Article{26}
We find that the number $+5$ is to be assigned to the collection
\[
\begin{array}{ll}
A\text{ for }p=&101,\ 109,\ 149,\ 181,\ 269\\
B\text{ for }p=&13,\ 17,\ 73,\ 97,\ 157,\ 193,\ 197,\ 233,\ 277,\ 293\\
C\text{ for }p=&29,\ 41,\ 61,\ 89,\ 229,\ 241,\ 281\\
D\text{ for }p=&37,\ 53,\ 113,\ 137,\ 173,\ 257
\end{array}
\]
Considering the values of the numbers $a,b$ corresponding to each $p$, we find the law here as readily as for the classification of the number $-3$. Namely, we arrive at the collection

\SourceLeaf{108}%
\[
\begin{array}{ll}
A,&\text{whenever }b\equiv0\ (\mathrm{mod.}\ 5)\\
B,&\text{whenever }b\equiv a\\
C,&\text{whenever }a\equiv0\\
D,&\text{whenever }b\equiv4a
\end{array}
\]

These rules evidently cover all cases, since for $b\equiv2a$ or $b\equiv3a\ (\mathrm{mod.}\ 5)$ we would have $aa+bb\equiv0$, Q. E. A., since by hypothesis $p$ is a prime number different from $5$.

\Article{27}
Likewise, inductive investigation applied to the numbers $-7$, $-11$, $+13$, $+17$, $-19$, $-23$ and carried far enough points to the following rules:
\begin{center}For the number $-7$\end{center}
\[
\begin{array}{c|l}
A&a\equiv0,\quad\text{or }b\equiv0\ (\mathrm{mod.}\ 7)\\
B&b\equiv4a,\ \text{or }b\equiv5a\\
C&b\equiv a,\quad\text{or }b\equiv6a\\
D&b\equiv2a,\ \text{or }b\equiv3a
\end{array}
\]

\begin{center}For the number $-11$.\end{center}
\[
\begin{array}{c|l}
A&b\equiv0,\ 5a,\ \text{or }6a\ (\mathrm{mod.}\ 11)\\
B&b\equiv a,\ 3a\text{ or }4a\\
C&a\equiv0,\ \text{or }b\equiv2a\text{ or }9a\\
D&b\equiv7a,\ 8a\text{ or }10a
\end{array}
\]

\begin{center}For the number $+13$.\end{center}
\[
\begin{array}{c|l}
A&b\equiv0,\ 4a,\ 9a\ (\mathrm{mod.}\ 13)\\
B&b\equiv6a,\ 11a,\ 12a\\
C&a\equiv0;\ b\equiv3a,\ 10a\\
D&b\equiv a,\ 2a,\ 7a
\end{array}
\]

\begin{center}For the number $+17$.\end{center}
\[
\begin{array}{c|l}
A&a\equiv0;\ b\equiv0,\ a,\ 16a\ (\mathrm{mod.}\ 17)\\
B&b\equiv2a,\ 6a,\ 8a,\ 14a\\
C&b\equiv5a,\ 7a,\ 10a,\ 12a\\
D&b\equiv3a,\ 9a,\ 11a,\ 15a
\end{array}
\]

\SourceLeaf{109}%
\begin{center}For the number $-19$.\end{center}
\[
\begin{array}{c|l}
A&b\equiv0,\ 2a,\ 5a,\ 14a,\ 17a\ (\mathrm{mod.}\ 19)\\
B&b\equiv3a,\ 7a,\ 11a,\ 13a,\ 18a\\
C&a\equiv0;\ b\equiv4a,\ 9a,\ 10a,\ 15a\\
D&b\equiv a,\ 6a,\ 8a,\ 12a,\ 16a
\end{array}
\]

\begin{center}For the number $-23$.\end{center}
\[
\begin{array}{c|l}
A&a\equiv0;\ b\equiv0,\ 7a,\ 10a,\ 13a,\ 16a\ (\mathrm{mod.}\ 23)\\
B&b\equiv2a,\ 3a,\ 4a,\ 11a,\ 15a,\ 17a\\
C&b\equiv a,\ 5a,\ 9a,\ 14a,\ 18a,\ 22a\\
D&b\equiv6a,\ 8a,\ 12a,\ 19a,\ 20a,\ 21a
\end{array}
\]

\Article{28}
The special theorems obtained in this way by inductive investigation are found to be confirmed as far as that investigation is continued, and reveal a particularly beautiful form of the criteria. But when they are compared with one another to draw general conclusions, the following observations immediately present themselves.

The criteria for deciding which collection the prime number $\pm q$ belongs to (taking the upper or lower sign according as $q$ has the form $4n+1$ or $4n+3$) depend on the forms of the numbers $a,b$ compared with one another with respect to the modulus $q$. Namely,

\Needspace{11\baselineskip}
I. Whenever $a\equiv0\ (\mathrm{mod.}\ q)$, $\pm q$ belongs to a definite collection, which is $A$ for $q=7,\ 17,\ 23$, and $C$ for $q=3,\ 11,\ 13,\ 19$. This suggests the conjecture that the former case holds generally whenever $q$ has the form $8n\pm1$, and the latter whenever $q$ has the form $8n\pm3$. Moreover, the collections $B$ and $D$ can already be excluded without inductive investigation when the value of $a$ is divisible by $q$, in which case $p\equiv bb\ (\mathrm{mod.}\ q)$, that is, $p$ is a quadratic residue of $q$; hence, by the fundamental theorem, $\pm q$ must be a quadratic residue of $p$.

II. When $a$ is not divisible by $q$, the criterion depends on the value of the expression $\frac ba\ (\mathrm{mod.}\ q)$. This expression does admit $q$ different values, namely $0,\ 1,\ 2,\ 3\dotsfour q-1$; but whenever $q$ has the form $4n+1$, the two
\SourceLeaf{110}%
values of the expression $\surd-1\ (\mathrm{mod.}\ q)$ must be excluded. They evidently cannot be values of the expression $\frac ba\ (\mathrm{mod.}\ q)$, since $p=aa+bb$ is always assumed to be a prime number different from $q$. The number of admissible values of the expression $\frac ba\ (\mathrm{mod.}\ q)$ is therefore $=q-2$ for $q\equiv1\ (\mathrm{mod.}\ 4)$, while it remains $=q$ for $q\equiv3\ (\mathrm{mod.}\ 4)$.

These values are now distributed among four classes: some, to be denoted without further specification by $\alpha$, correspond to the collection $A$; others, denoted by $\beta$, to the collection $B$; others, $\gamma$, to the collection $C$; and finally the rest, $\delta$, to the collection $D$. Thus $\pm q$ is to be assigned to the collection $A,B,C,D$ according as $b\equiv\alpha a$, $b\equiv\beta a$, $b\equiv\gamma a$, $b\equiv\delta a\ (\mathrm{mod.}\ q)$.

But the \emph{law} of this distribution seems more deeply hidden, although some general features are readily noticed. The number of values in three of the classes is found to be the same, namely $=\frac14(q-1)$ or $\frac14(q+1)$, while in one (specifically, the one corresponding to the collection with the criterion $a\equiv0$) it is smaller by one. Thus the total number of different criteria for each collection is the same, namely $=\frac14(q-1)$ or $\frac14(q+1)$. We further observe that $0$ is always found in the first class (among the $\alpha$), and that the complements of the numbers $\alpha,\beta,\gamma,\delta$ to $q$, namely $q-\alpha$, $q-\beta$, $q-\gamma$, $q-\delta$, are found in the first, fourth, third and second classes, respectively. Finally, we see that the values of the expressions $\frac1\alpha$, $\frac1\beta$, $\frac1\gamma$, $\frac1\delta\ (\mathrm{mod.}\ q)$ belong to the first, fourth, third and second classes whenever the criterion $a\equiv0$ corresponds to the collection $A$, but to the third, second, first and fourth classes, respectively, whenever the criterion $a\equiv0$ refers to the collection $C$. But this is about the limit of what can be reached by inductive investigation, unless we rather boldly presume to anticipate what will later be drawn from the genuine sources.
